% Generated by paper-covert from output/manuscript/sections_revised/04_study_area_and_data.md

\section{Study Area and Data}\label{study-area-and-data}

\subsection{Study area}\label{study-area}

The completed national map covers continent-wide Australia, comprising 99,465 tiles (98.9\% of the national extent that National Land Use Map, NLUM, probability surfaces indicate carries meaningful agricultural probability for one of the three target crop groups) for the 2024 growing season. A 100 km $\times$ 100 km block in the Riverina region of New South Wales, a major mixed-cropping district, was used for an intensively validated nine-season (2017-2025) regional test of the identical pipeline, reported separately from the national result because only the regional extent currently has multi-year coverage (Section 3.2; Section 7, Limitation 1). Fig. 2 shows the national tile coverage, the Riverina regional block, and the SA2 regions used for validation (Section 3.3).

\begin{figure}[htbp]
\centering
\includegraphics[width=0.85\textwidth]{figures/Fig02_study_area}
\caption{Study area. National extent of 2024 mapped tiles (grey; 99,465 tile centres), the 100km $\times$ 100km Riverina, NSW block used for the nine-year (2017-2025) regional test (red outline), and the SA2s used for ABS composition validation (coloured by state; EPSG:3577). 133 SA2s shown --- the $\geq$50\%-covered subset re-derived from the published per-SA2 cover table.}
\label{fig:2}
\end{figure}

\subsection{Temporal scope}\label{temporal-scope}

A single growing season, 2024, is complete and validated at national scale. Nine seasons (2017-2025) are complete at the 100 km regional scale. A national extension to the remaining eight seasons (2017-2023, 2025) is planned but has not yet been run at the time of this draft {[}PENDING: national 2017-2025 run, E8{]}; all results and figures in this paper that describe multi-year behaviour are explicitly regional, not national, and are reported as suggestive of but not proof of national multi-year generalization (Section 6.6).

\subsection{Data sources}\label{data-sources}

Table 1 summarizes the datasets used. Sentinel-2 surface-reflectance imagery was accessed via Digital Earth Australia's Analysis Ready Data product (\texttt{ga\_s2am/bm/cm\_ard\_3}) through the National Computational Infrastructure's datacube, at 10 m resolution, nationally, 2017-2025. Ground truth for training and validation is drawn from the Grains Research and Development Corporation (GRDC) National Variety Trial (NVT) network: field-verified, point-located records of sown crop per trial site, spanning 2017-2025 nationally. This dataset is covered by a data-use agreement and cannot be released or described at the site level (coordinates, trial identifiers, or yields); all reporting in this paper is aggregate. Independent validation uses ABS sown-area statistics, aggregated to the SA2 level, which have no connection to the training pipeline; ABARES production statistics provide a secondary cross-check on ABS itself, and the two agencies' median 6.8\% mutual disagreement on national sown area is the accuracy floor this paper measures its own map against (Section 5.2). NLUM v7 250 m per-commodity agricultural probability surfaces (2020-21) were used during national-scale inference to restrict segmentation and classification to tiles with non-trivial probability for one of the three target crop groups, and were separately used post hoc as an independent spatial-comparison product (Section 5.5); this is a different use of the same dataset and is described separately in each place it appears. Paddock boundaries were delineated with the Segment Anything Model (\texttt{vit\_h} checkpoint, stock parameters; \citealp{kirillov2023sam}) via the SAMGeo/PaddockTS wrapper; re-tuning SAM's segmentation parameters was investigated as a possible fix for a subset of poorly segmented polygons and ruled out as not cost-effective relative to the presence-gate and area-mask approach actually shipped.

\textbf{Table 1.} Dataset and pipeline summary.

\begin{longtable}[]{@{}
  >{\raggedright\arraybackslash}p{(\linewidth - 6\tabcolsep) * \real{0.2500}}
  >{\raggedright\arraybackslash}p{(\linewidth - 6\tabcolsep) * \real{0.2500}}
  >{\raggedright\arraybackslash}p{(\linewidth - 6\tabcolsep) * \real{0.2500}}
  >{\raggedright\arraybackslash}p{(\linewidth - 6\tabcolsep) * \real{0.2500}}@{}}
\toprule\noalign{}
\begin{minipage}[b]{\linewidth}\raggedright
Dataset
\end{minipage} & \begin{minipage}[b]{\linewidth}\raggedright
Role
\end{minipage} & \begin{minipage}[b]{\linewidth}\raggedright
Resolution / extent
\end{minipage} & \begin{minipage}[b]{\linewidth}\raggedright
Access
\end{minipage} \\
\midrule\noalign{}
\endhead
\bottomrule\noalign{}
\endlastfoot
Sentinel-2 (DEA ARD) & Primary imagery; day-of-year-binned spectral indices & 10 m, national, 2017-2025 & Open \\
GRDC / NVT trial records & Training and validation ground truth & Point-in-paddock, national, 2017-2025 & Data-use agreement; not releasable; aggregate reporting only \\
ABS sown-area statistics & Independent validation reference & SA2-level, 133 SA2s, national & Public \\
ABARES production statistics & Cross-check on ABS (6.8\% median mutual disagreement) & State / national & Public \\
NLUM v7 250 m Ag. Probability Surfaces & (a) National-run coverage mask; (b) independent spatial-comparison product & 250 m, national & Public \\
WorldCereal 2021 v100 & Independent spatial-comparison product & 10 m, national AU extent & Public (CC-BY 4.0) \\
Segment Anything Model (SAM) & Paddock-scale segmentation, stock parameters & 10 m-derived polygons & Open weights \\
\end{longtable}


\subsection{Preprocessing}\label{preprocessing}

Segmented polygons were retained only within a 5-300 ha area range with a compactness ratio (perimeter / sqrt(area)) of at most 8 --- a free, non-spectral presence signal derived purely from segmentation geometry, which retains 77.9\% of known crop paddocks in validation while rejecting 65\% of pasture polygons, at a cost of excluding 22\% of real crop paddocks that fall outside the size/shape mask. Fifty-one paddock-median spectral features were computed per retained polygon, binned by day-of-year across the growing season, from three indices: NDVI, a canola-specific flowering index following the general approach of \citet{ashourloo2019canola}, and a chlorophyll-fluorescence-related index. NVT trial points were matched to their enclosing SAM polygon; because 69.7\% of training trials share a segmented paddock with another NVT trial (31.8\% with a different crop, reflecting how trial networks co-locate multiple varieties), the original nine-species label scheme was collapsed to three broad groups, resolving 69\% of these label conflicts, and the remainder were resolved by hand review (Section 4.2).
